\chapter{Complexity — Code Parser}
%%% generated by the blueprint pipeline from TCSlib.Complexity.TuringMachine.CodeParser
%%%%%%%%%%%%%%%%%%%%%%%%%%%%%%%%%%%%%%%%%%%%%%%%%%%%%%%%%%%%%%%%%%%%%%%%%%%%%%%
\section{Overview}

This module documents the parser and decoder layer for the fixed bit-string
serialization of coded Turing machines (one-work-tape machines over the binary
alphabet, bundled with their state count).  It provides field readers for each
component of the serialization grammar, a total decoder mapping every malformed
string to a fixed fallback machine together with its padded round-trip law,
soundness of the parser (acceptance characterizes a canonical serialization
followed by all-1 padding), and an erased suffix scanner underlying a canonizer
that computes the serialization of the machine a string denotes.  Throughout,
bit strings are written over $\{0,1\}$, with $1$ for true and $0$ for false, and
a \emph{reader} maps a bit string either to failure or to a decoded value paired
with the unconsumed suffix.

%%%%%%%%%%%%%%%%%%%%%%%%%%%%%%%%%%%%%%%%%%%%%%%%%%%%%%%%%%%%%%%%%%%%%%%%%%%%%%%
\section{Declarations}

\begin{definition}[Unary numeral reader]
\lean{Turing.codeReadUnary}
\label{Turing.codeReadUnary}
\leanok
A reader for unary numerals: on an input beginning with $n$ ones followed by a
zero it succeeds, returning the natural number $n$ together with the suffix
after that zero; it fails if the input is exhausted before a zero is found.
\end{definition}

\begin{lemma}[Unary reader is a prefix inverse]
\lean{Turing.codeReadUnary_append}
\label{Turing.codeReadUnary_append}
\leanok
\uses{Turing.codeReadUnary}
\difficulty{3}
For every $n \in \mathbb{N}$ and every bit string $x$, the unary numeral reader
applied to the string consisting of $n$ ones, then a zero, then $x$, succeeds
and returns the pair $(n, x)$; in particular the suffix $x$ is left untouched.
\end{lemma}

\begin{definition}[Range-checked unary index reader]
\lean{Turing.codeReadFin}
\label{Turing.codeReadFin}
\leanok
\uses{Turing.codeReadUnary}
Given a bound $n$, the reader that decodes a unary numeral (ones terminated by
a zero) from a bit string and succeeds only when the decoded value $i$
satisfies $i < n$, in which case it returns $i$ as an element of
$\{0, \dots, n-1\}$ together with the unconsumed suffix.
\end{definition}

\begin{definition}[Head-movement dictionary reader]
\lean{Turing.codeReadSign}
\label{Turing.codeReadSign}
\leanok
A reader decoding a fixed two-bit dictionary for a head movement: a leading
$11$ denotes moving left, $00$ staying put, and $10$ moving right, in each case
returning the movement together with the remaining suffix; a leading $01$, or
an input of fewer than two bits, is rejected.
\end{definition}

\begin{definition}[Optional-output dictionary reader]
\lean{Turing.codeReadOutput}
\label{Turing.codeReadOutput}
\leanok
A reader decoding a fixed two-bit dictionary for an optional output bit: a
leading $00$ denotes producing no output, $10$ outputting the bit $0$, and $11$
outputting the bit $1$; a leading $01$, or an input of fewer than two bits, is
rejected.  On success the decoded value is returned with the remaining suffix.
\end{definition}

\begin{definition}[Optional-write dictionary reader]
\lean{Turing.codeReadWrite}
\label{Turing.codeReadWrite}
\leanok
A reader decoding a fixed two-bit dictionary for an optional work-tape write,
whose possible values are: write nothing, write the blank symbol, write the bit
$0$, or write the bit $1$, denoted by the leading pairs $00$, $01$, $10$, $11$
respectively.  Inputs of fewer than two bits are rejected; on success the
decoded value is returned with the remaining suffix.
\end{definition}

\begin{definition}[Successor-state field reader]
\lean{Turing.codeReadState}
\label{Turing.codeReadState}
\leanok
\uses{Turing.codeReadFin}
Given a bound $n$, a reader for a successor-state field: a leading $0$ denotes
halting (no successor state) and consumes only that bit, while a leading $1$ is
followed by a unary index that must be smaller than $n$ and is returned as the
live successor state, in each case together with the unconsumed suffix.
\end{definition}

\begin{definition}[Transition-record reader]
\lean{Turing.codeReadAction}
\label{Turing.codeReadAction}
\leanok
\uses{Turing.Action, Turing.codeReadOutput, Turing.codeReadSign, Turing.codeReadState, Turing.codeReadWrite}
Given a state-count parameter $n$, a reader for one transition record of a
one-work-tape machine over the binary alphabet with $n + 1$ states.  It reads,
in order: an input-head movement, an optional work-tape write, a work-head
movement, an optional output bit, and a successor-state field over $n + 1$
states, failing as soon as any of the five fields fails, and assembles the
results into a single transition record together with the unconsumed suffix.
\end{definition}

\begin{definition}[Symbol-indexed row reader]
\lean{Turing.codeReadSymbols}
\label{Turing.codeReadSymbols}
\leanok
Given a field reader producing values in a type $A$, the reader that decodes
three consecutive $A$-entries and returns the function assigning them to the
three tape symbols blank, $0$, $1$ in that fixed order, together with the
unconsumed suffix; it fails as soon as any of the three reads fails.
\end{definition}

\begin{definition}[Fixed-length vector reader]
\lean{Turing.codeReadVec}
\label{Turing.codeReadVec}
\leanok
Given a field reader producing values in a type $A$ and a length $n$, the
reader that decodes $n$ consecutive $A$-entries and returns them as a vector
indexed by $\{0, \dots, n-1\}$, together with the unconsumed suffix; it fails
as soon as any entry fails.
\end{definition}

\begin{definition}[Little-endian binary value of a bit string]
\lean{Turing.codeBitsNat}
\label{Turing.codeBitsNat}
\leanok
The natural number denoted by a bit string read in least-significant-bit-first
binary, so the string $b_0 b_1 \cdots b_{k-1}$ denotes $\sum_{j<k} b_j 2^j$.
The function is total: it also assigns a value to noncanonical strings (for
instance those with trailing zeros); canonicity is checked separately by the
parser.
\end{definition}

\begin{definition}[Fallback machine]
\lean{Turing.codeFallback}
\label{Turing.codeFallback}
\leanok
\uses{Turing.CodeTM}
The fixed fallback coded machine: a machine with a single state that
immediately halts without moving any head, writing to the work tape, or
producing output.
\end{definition}

\begin{definition}[Machine parser]
\lean{Turing.codeParse}
\label{Turing.codeParse}
\leanok
\uses{Turing.CodeTM, Turing.codeBitsNat, Turing.codeReadAction, Turing.codeReadFin, Turing.codeReadSymbols, Turing.codeReadVec, Turing.pairDecode}
The partial parser for the fixed serialization of coded machines.  On a bit
string it: reads a doubled-bit header (each data bit repeated twice, terminated
by the separator $01$) and requires the recovered word to be the canonical
little-endian binary representation of its value $n$; rejects unless the
remaining input has length at least $81(n+1)$; reads a unary initial state
smaller than $n + 1$; reads the full transition table, namely for each of the
$n + 1$ states a three-by-three grid of transition records indexed by the input
and work symbols in blank/$0$/$1$ order; and finally requires every remaining
bit to be $1$.  On success it returns the coded machine with $n + 1$ states
assembled from these data.
\end{definition}

\begin{definition}[Total decoder]
\lean{Turing.codeDecode}
\label{Turing.codeDecode}
\leanok
\uses{Turing.CodeTM, Turing.codeFallback, Turing.codeParse}
The total decoding function from bit strings to coded machines: apply the
machine parser, and return the fixed fallback machine (one state, immediately
halting and silent) whenever parsing fails.
\end{definition}

\begin{lemma}[Bounded-index reader is a prefix inverse]
\lean{Turing.codeReadFin_append}
\label{Turing.codeReadFin_append}
\leanok
\uses{Turing.codeReadFin, Turing.codeReadUnary_append, Turing.unaryFin}
\difficulty{1}
For a bound $n$, an index $i < n$, and any bit string $x$: the range-checked
unary reader with bound $n$, applied to the self-delimiting unary code of $i$
($i$ ones followed by a zero) concatenated with $x$, succeeds and returns
$(i, x)$.
\end{lemma}

\begin{lemma}[Movement reader is a prefix inverse]
\lean{Turing.codeReadSign_append}
\label{Turing.codeReadSign_append}
\leanok
\uses{Turing.codeReadSign, Turing.signBits}
\difficulty{2}
For every head movement $s$ (left, stay, or right) and every bit string $x$,
the head-movement reader applied to the two-bit dictionary code of $s$ ($11$,
$00$, or $10$ respectively) followed by $x$ succeeds and returns $(s, x)$.
\end{lemma}

\begin{lemma}[Output reader is a prefix inverse]
\lean{Turing.codeReadOutput_append}
\label{Turing.codeReadOutput_append}
\leanok
\uses{Turing.codeReadOutput, Turing.optBoolBits}
\difficulty{2}
For every optional output bit $b$ (no output, output $0$, or output $1$) and
every bit string $x$, the optional-output reader applied to the two-bit
dictionary code of $b$ ($00$, $10$, or $11$ respectively) followed by $x$
succeeds and returns $(b, x)$.
\end{lemma}

\begin{lemma}[Write reader is a prefix inverse]
\lean{Turing.codeReadWrite_append}
\label{Turing.codeReadWrite_append}
\leanok
\uses{Turing.codeReadWrite, Turing.optOptBoolBits}
\difficulty{2}
For every optional work-tape write $b$ (write nothing, write the blank symbol,
write $0$, or write $1$) and every bit string $x$, the optional-write reader
applied to the two-bit dictionary code of $b$ ($00$, $01$, $10$, or $11$
respectively) followed by $x$ succeeds and returns $(b, x)$.
\end{lemma}

\begin{lemma}[Successor reader is a prefix inverse]
\lean{Turing.codeReadState_append}
\label{Turing.codeReadState_append}
\leanok
\uses{Turing.codeReadFin_append, Turing.codeReadState, Turing.optStateBits}
\difficulty{2}
For a bound $n$, a successor-state field $s$ (either halt or a live state
$i < n$), and any bit string $x$: the successor-state reader with bound $n$,
applied to the encoding of $s$ (a single $0$ for halt; a $1$ followed by the
unary code of $i$ for a live state) followed by $x$, succeeds and returns
$(s, x)$.
\end{lemma}

\begin{lemma}[Record reader is a prefix inverse]
\lean{Turing.codeReadAction_append}
\label{Turing.codeReadAction_append}
\leanok
\uses{Turing.Action, Turing.actionBits, Turing.codeReadAction, Turing.codeReadOutput_append, Turing.codeReadSign_append, Turing.codeReadState_append, Turing.codeReadWrite_append}
\difficulty{4}
For every transition record $a$ of a one-work-tape binary machine with $n + 1$
states and every bit string $x$, the transition-record reader applied to the
serialization of $a$ (the dictionary codes of its input-head movement,
work-tape write, work-head movement, and output, followed by its
successor-state field) concatenated with $x$ succeeds and returns $(a, x)$.
\end{lemma}

\begin{lemma}[Row reader is a prefix inverse]
\lean{Turing.codeReadSymbols_append}
\label{Turing.codeReadSymbols_append}
\leanok
\uses{Turing.codeReadSymbols}
\difficulty{3}
Let a field reader and a field writer for values in a type $A$ be given such
that reading from any written value followed by any suffix returns exactly that
value and suffix.  Then for every function $f$ from the three tape symbols
(blank, $0$, $1$) to $A$ and every bit string $x$, the three-symbol row reader
applied to the concatenation of the written values $f(\mathrm{blank})$, $f(0)$,
$f(1)$ in that order, followed by $x$, succeeds and returns $(f, x)$.
\end{lemma}

\begin{lemma}[Vector reader is a prefix inverse]
\lean{Turing.codeReadVec_append}
\label{Turing.codeReadVec_append}
\leanok
\uses{Turing.codeReadVec}
\difficulty{4}
Let a field reader and a field writer for values in a type $A$ be given such
that reading from any written value followed by any suffix returns exactly that
value and suffix.  Then for every $n$, every vector $f$ indexed by
$\{0, \dots, n-1\}$, and every bit string $x$, the length-$n$ vector reader
applied to the concatenation of the written entries $f(0), \dots, f(n-1)$ in
order, followed by $x$, succeeds and returns $(f, x)$.
\end{lemma}

\begin{lemma}[Binary value inverts canonical bits]
\lean{Turing.codeBitsNat_bits}
\label{Turing.codeBitsNat_bits}
\leanok
\uses{Turing.codeBitsNat}
\difficulty{3}
For every natural number $n$, interpreting the canonical
least-significant-bit-first binary representation of $n$ (the empty string for
$n = 0$) as a binary numeral yields $n$ again.
\end{lemma}

\begin{lemma}[Lower bound on record length]
\lean{Turing.codeAction_length}
\label{Turing.codeAction_length}
\leanok
\uses{Turing.Action, Turing.actionBits, Turing.optBoolBits, Turing.optOptBoolBits, Turing.optStateBits, Turing.signBits, Turing.unaryFin}
\difficulty{4}
Every serialized transition record of a one-work-tape binary machine with
$n + 1$ states has length at least $9$: its four two-bit dictionary fields
contribute eight bits, and its successor-state field is nonempty.
\end{lemma}

\begin{lemma}[Length bound for concatenated words]
\lean{Turing.codeFlatMap_length}
\label{Turing.codeFlatMap_length}
\leanok
\difficulty{3}
If every element of a finite list maps under a function $f$ to a bit string of
length at least $c$, then the concatenation of all these bit strings, in list
order, has length at least $c$ times the length of the list.
\end{lemma}

\begin{lemma}[Lower bound on table length]
\lean{Turing.codeTable_length}
\label{Turing.codeTable_length}
\leanok
\uses{Turing.CodeTM, Turing.actionBits, Turing.codeAction_length, Turing.codeFlatMap_length}
\difficulty{4}
For every coded machine with state-count parameter $n$ (so $n + 1$ states), the
concatenation of the serialized transition records over all states and all nine
input/work symbol pairs has length at least $81(n+1)$.
\end{lemma}

\begin{lemma}[Table reader is a prefix inverse]
\lean{Turing.codeReadTable_append}
\label{Turing.codeReadTable_append}
\leanok
\uses{Turing.CodeTM, Turing.actionBits, Turing.codeReadAction, Turing.codeReadAction_append, Turing.codeReadSymbols, Turing.codeReadSymbols_append, Turing.codeReadVec, Turing.codeReadVec_append}
\difficulty{2}
Let $M$ be a coded machine with state-count parameter $n$.  Applying the table
reader --- the length-$(n+1)$ vector reader whose entries are three-by-three
grids of transition records indexed by the input and work symbols --- to the
concatenated serializations of all of $M$'s transition records (states in
order, input and work symbols each in blank/$0$/$1$ order) followed by any bit
string $x$ succeeds, returning the function that assigns to each state, input
symbol, and work symbol the corresponding transition of $M$, together with $x$
untouched.
\end{lemma}

\begin{lemma}[Parser round-trip under padding]
\lean{Turing.codeParse_serialize_pad}
\label{Turing.codeParse_serialize_pad}
\leanok
\uses{Turing.CodeTM, Turing.CodeTM.serialize, Turing.codeBitsNat_bits, Turing.codeParse, Turing.codeReadFin_append, Turing.codeReadTable_append, Turing.codeTable_length, Turing.pairDecode_pairEncode, Turing.pairEncode}
\difficulty{5}
For every coded machine $M$ and every $m \in \mathbb{N}$, the machine parser
applied to the canonical serialization of $M$ followed by $m$ additional
$1$-bits succeeds and returns exactly $M$.
\end{lemma}

\begin{lemma}[Decoder round-trip under padding]
\lean{Turing.codeDecode_serialize_pad}
\label{Turing.codeDecode_serialize_pad}
\leanok
\uses{Turing.CodeTM, Turing.CodeTM.serialize, Turing.codeDecode, Turing.codeParse_serialize_pad}
\difficulty{1}
For every coded machine $M$ and every $m \in \mathbb{N}$, the total decoder
maps the canonical serialization of $M$ followed by $m$ additional $1$-bits
back to $M$.
\end{lemma}

\begin{lemma}[Soundness of the unary reader]
\lean{Turing.codeReadUnary_sound}
\label{Turing.codeReadUnary_sound}
\leanok
\uses{Turing.codeReadUnary}
\difficulty{4}
If the unary numeral reader succeeds on a bit string $x$, returning the value
$n$ and the suffix $r$, then $x$ is exactly $n$ ones, followed by a zero,
followed by $r$.
\end{lemma}

\begin{lemma}[Soundness of the bounded-index reader]
\lean{Turing.codeReadFin_sound}
\label{Turing.codeReadFin_sound}
\leanok
\uses{Turing.codeReadFin, Turing.codeReadUnary_sound, Turing.unaryFin}
\difficulty{3}
If the range-checked unary reader with bound $n$ succeeds on a bit string $x$,
returning an index $i < n$ and a suffix $r$, then $x$ is exactly the
self-delimiting unary code of $i$ ($i$ ones and a closing zero) followed by
$r$.
\end{lemma}

\begin{lemma}[Soundness of the doubled-bit header parser]
\lean{Turing.codePairDecode_sound}
\label{Turing.codePairDecode_sound}
\leanok
\uses{Turing.pairDecode, Turing.pairEncode}
\difficulty{5}
If the doubled-bit header parser succeeds on a bit string $x$, returning a data
word $a$ and a suffix $r$, then $x$ is exactly the doubled-bit encoding of $a$
--- each bit of $a$ repeated twice, then the separator $01$ --- followed by
$r$.  This includes the case of an empty data word.
\end{lemma}

\begin{lemma}[Soundness of the movement reader]
\lean{Turing.codeReadSign_sound}
\label{Turing.codeReadSign_sound}
\leanok
\uses{Turing.codeReadSign, Turing.signBits}
\difficulty{3}
If the head-movement reader succeeds on a bit string $x$, returning a movement
$s$ and a suffix $r$, then $x$ is exactly the two-bit dictionary code of $s$
followed by $r$.
\end{lemma}

\begin{lemma}[Soundness of the output reader]
\lean{Turing.codeReadOutput_sound}
\label{Turing.codeReadOutput_sound}
\leanok
\uses{Turing.codeReadOutput, Turing.codeReadWrite_sound, Turing.optBoolBits}
\difficulty{3}
If the optional-output reader succeeds on a bit string $x$, returning a value
$b$ and a suffix $r$, then $x$ is exactly the two-bit dictionary code of $b$
followed by $r$.
\end{lemma}

\begin{lemma}[Soundness of the write reader]
\lean{Turing.codeReadWrite_sound}
\label{Turing.codeReadWrite_sound}
\leanok
\uses{Turing.codeReadWrite, Turing.optOptBoolBits}
\difficulty{3}
If the optional-write reader succeeds on a bit string $x$, returning a value
$b$ and a suffix $r$, then $x$ is exactly the two-bit dictionary code of $b$
followed by $r$.
\end{lemma}

\begin{lemma}[Soundness of the successor reader]
\lean{Turing.codeReadState_sound}
\label{Turing.codeReadState_sound}
\leanok
\uses{Turing.codeReadFin, Turing.codeReadFin_sound, Turing.codeReadState, Turing.optStateBits}
\difficulty{4}
If the successor-state reader with bound $n$ succeeds on a bit string $x$,
returning a field $s$ and a suffix $r$, then $x$ is exactly the encoding of $s$
(a single $0$ for halt, or a $1$ followed by the unary code of a live state
below $n$) followed by $r$.
\end{lemma}

\begin{lemma}[Soundness of the record reader]
\lean{Turing.codeReadAction_sound}
\label{Turing.codeReadAction_sound}
\leanok
\uses{Turing.Action, Turing.actionBits, Turing.codeReadAction, Turing.codeReadOutput_sound, Turing.codeReadSign_sound, Turing.codeReadState_sound, Turing.codeReadWrite_sound}
\difficulty{4}
If the transition-record reader with state-count parameter $n$ succeeds on a
bit string $x$, returning a record $a$ and a suffix $r$, then $x$ is exactly
the serialization of $a$ --- the codes of its five fields concatenated in order
--- followed by $r$.
\end{lemma}

\begin{lemma}[Soundness of the row reader]
\lean{Turing.codeReadSymbols_sound}
\label{Turing.codeReadSymbols_sound}
\leanok
\uses{Turing.codeReadSymbols, Turing.codeScan}
\difficulty{3}
Let a field reader and a field writer for values in a type $A$ be such that
every successful read on a string exhibits that string as the written value
followed by the returned suffix.  Then whenever the three-symbol row reader
succeeds on a bit string $x$, returning a row $f$ and a suffix $r$, the string
$x$ is exactly the concatenation of the written values $f(\mathrm{blank})$,
$f(0)$, $f(1)$ followed by $r$.
\end{lemma}

\begin{lemma}[Soundness of the vector reader]
\lean{Turing.codeReadVec_sound}
\label{Turing.codeReadVec_sound}
\leanok
\uses{Turing.codeReadVec}
\difficulty{4}
Let a field reader and a field writer for values in a type $A$ be such that
every successful read on a string exhibits that string as the written value
followed by the returned suffix.  Then whenever the length-$n$ vector reader
succeeds on a bit string $x$, returning a vector $f$ and a suffix $r$, the
string $x$ is exactly the in-order concatenation of the written entries
$f(0), \dots, f(n-1)$ followed by $r$.
\end{lemma}

\begin{lemma}[All-true strings are replicates]
\lean{Turing.codeAllTrue_eq}
\label{Turing.codeAllTrue_eq}
\leanok
\difficulty{3}
A bit string all of whose bits are $1$ is equal to the constant-$1$ string of
its own length.
\end{lemma}

\begin{lemma}[Soundness of the machine parser]
\lean{Turing.codeParse_sound}
\label{Turing.codeParse_sound}
\leanok
\uses{Turing.CodeTM, Turing.CodeTM.serialize, Turing.codeAllTrue_eq, Turing.codeBitsNat, Turing.codePairDecode_sound, Turing.codeParse, Turing.codeReadAction_sound, Turing.codeReadFin_sound, Turing.codeReadSymbols_sound, Turing.codeReadVec_sound, Turing.codeScan_full, Turing.pairEncode}
\difficulty{5}
If the machine parser accepts a bit string $x$, returning a coded machine $M$,
then $x$ is the canonical serialization of $M$ followed by some number of
$1$-bits.
\end{lemma}

\begin{definition}[Two-bit skipper]
\lean{Turing.codeSkipPair}
\label{Turing.codeSkipPair}
\leanok
Given a validity predicate on pairs of bits, the reader that consumes exactly
two bits from its input, succeeds precisely when the consumed pair satisfies
the predicate, and returns only the unconsumed suffix.  Inputs of fewer than
two bits are rejected.
\end{definition}

\begin{definition}[Unary index skipper]
\lean{Turing.codeSkipFin}
\label{Turing.codeSkipFin}
\leanok
\uses{Turing.codeReadUnary}
Given a bound $n$, the reader that consumes a unary numeral (ones terminated by
a zero), succeeds precisely when its value is smaller than $n$, and returns
only the unconsumed suffix.
\end{definition}

\begin{definition}[Successor-field skipper]
\lean{Turing.codeSkipState}
\label{Turing.codeSkipState}
\leanok
\uses{Turing.codeSkipFin}
Given a bound $n$, the reader that consumes a successor-state field --- either
a single $0$ (halt), or a $1$ followed by a unary index that must be smaller
than $n$ --- and returns only the unconsumed suffix.
\end{definition}

\begin{definition}[Transition-record skipper]
\lean{Turing.codeSkipAction}
\label{Turing.codeSkipAction}
\leanok
\uses{Turing.codeSkipPair, Turing.codeSkipState}
Given a state-count parameter $n$, the reader that consumes one five-field
transition record: four two-bit dictionary fields (input-head movement,
work-tape write, work-head movement, output), each validated against its
dictionary, followed by a successor-state field over $n + 1$ states.  It
returns only the unconsumed suffix.
\end{definition}

\begin{definition}[Iterated skipper]
\lean{Turing.codeSkipRepeat}
\label{Turing.codeSkipRepeat}
\leanok
Given a skipping reader mapping a bit string to an optional suffix, and a count
$n$, the reader that applies the skipper $n$ times in sequence, threading the
suffix through, and returns the final suffix; it fails as soon as any
iteration fails.
\end{definition}

\begin{lemma}[Erasure law for the bounded-index reader]
\lean{Turing.codeEraseFin}
\label{Turing.codeEraseFin}
\leanok
\uses{Turing.codeReadFin, Turing.codeSkipFin}
\difficulty{3}
For every bound $n$ and bit string $x$: discarding the parsed index from the
result of the range-checked unary reader (keeping only the suffix, or failure)
yields exactly the result of the unary index skipper with bound $n$ on $x$.
\end{lemma}

\begin{lemma}[Erasure law for the movement reader]
\lean{Turing.codeEraseSign}
\label{Turing.codeEraseSign}
\leanok
\uses{Turing.codeReadSign, Turing.codeSkipPair}
\difficulty{2}
For every bit string $x$: discarding the parsed movement from the head-movement
reader yields exactly the two-bit skipper whose validity predicate accepts a
pair $(a, b)$ precisely when $a$ is $1$ or $b$ is $0$ (that is, it rejects only
the pair $01$).
\end{lemma}

\begin{lemma}[Erasure law for the output reader]
\lean{Turing.codeEraseOutput}
\label{Turing.codeEraseOutput}
\leanok
\uses{Turing.codeReadOutput, Turing.codeSkipPair}
\difficulty{2}
For every bit string $x$: discarding the parsed value from the optional-output
reader yields exactly the two-bit skipper whose validity predicate rejects only
the pair $01$.
\end{lemma}

\begin{lemma}[Erasure law for the write reader]
\lean{Turing.codeEraseWrite}
\label{Turing.codeEraseWrite}
\leanok
\uses{Turing.codeReadWrite, Turing.codeSkipPair}
\difficulty{2}
For every bit string $x$: discarding the parsed value from the optional-write
reader yields exactly the two-bit skipper that accepts every pair of bits.
\end{lemma}

\begin{lemma}[Erasure law for the successor reader]
\lean{Turing.codeEraseState}
\label{Turing.codeEraseState}
\leanok
\uses{Turing.codeEraseFin, Turing.codeReadState, Turing.codeSkipState}
\difficulty{2}
For every bound $n$ and bit string $x$: discarding the parsed field from the
successor-state reader with bound $n$ yields exactly the successor-field
skipper with bound $n$.
\end{lemma}

\begin{lemma}[Bind commutes with value erasure]
\lean{Turing.codeErase_bind}
\label{Turing.codeErase_bind}
\leanok
\difficulty{1}
For an optional parse result $r$ --- either failure or a pair of a value and a
suffix --- and a partial continuation $f$ on suffixes: binding $r$ through the
function applying $f$ to the suffix component equals first discarding the value
component of $r$ and then binding through $f$.
\end{lemma}

\begin{lemma}[Erasure law for the record reader]
\lean{Turing.codeEraseAction}
\label{Turing.codeEraseAction}
\leanok
\uses{Turing.codeEraseOutput, Turing.codeEraseSign, Turing.codeEraseState, Turing.codeEraseWrite, Turing.codeErase_bind, Turing.codeReadAction, Turing.codeSkipAction}
\difficulty{4}
For every state-count parameter $n$ and bit string $x$: discarding the parsed
record from the transition-record reader with parameter $n$ yields exactly the
transition-record skipper with parameter $n$, which validates the four
dictionary fields and the successor field without constructing the record.
\end{lemma}

\begin{lemma}[Erasure law for the row reader]
\lean{Turing.codeEraseSymbols}
\label{Turing.codeEraseSymbols}
\leanok
\uses{Turing.codeReadSymbols, Turing.codeSkipRepeat}
\difficulty{1}
For every field reader with values in a type $A$ and every bit string $x$:
discarding the parsed row from the three-symbol row reader equals iterating the
value-discarding version of the field reader three times.
\end{lemma}

\begin{lemma}[Erasure law for the vector reader]
\lean{Turing.codeEraseVec}
\label{Turing.codeEraseVec}
\leanok
\uses{Turing.codeErase_bind, Turing.codeReadVec, Turing.codeSkipRepeat}
\difficulty{4}
For every field reader with values in a type $A$, every length $n$, and every
bit string $x$: discarding the parsed vector from the length-$n$ vector reader
equals iterating the value-discarding version of the field reader $n$ times.
\end{lemma}

\begin{definition}[Suffix-retaining machine parser]
\lean{Turing.codeParseFull}
\label{Turing.codeParseFull}
\leanok
\uses{Turing.CodeTM, Turing.codeBitsNat, Turing.codeReadAction, Turing.codeReadFin, Turing.codeReadSymbols, Turing.codeReadVec, Turing.pairDecode}
A variant of the machine parser following exactly the same grammar ---
doubled-bit canonical state count, minimum-length guard, unary initial state,
full transition table, and an all-$1$ remainder --- but returning the decoded
coded machine together with the unconsumed all-$1$ suffix, rather than the
machine alone.
\end{definition}

\begin{definition}[Erased suffix scanner]
\lean{Turing.codeScan}
\label{Turing.codeScan}
\leanok
\uses{Turing.codeBitsNat, Turing.codeSkipAction, Turing.codeSkipFin, Turing.codeSkipRepeat, Turing.pairDecode}
The erased parser: it follows the same grammar as the machine parser ---
doubled-bit canonical state count $n$, minimum-length guard $81(n+1)$, a
skipped unary initial state below $n + 1$, then $(n+1) \cdot 3 \cdot 3$ skipped
transition records, and an all-$1$ remainder --- but validates the fields
without constructing any machine, returning only the unconsumed all-$1$
suffix.
\end{definition}

\begin{lemma}[Parser as projection of the suffix-retaining parser]
\lean{Turing.codeParse_full}
\label{Turing.codeParse_full}
\leanok
\uses{Turing.codeParse, Turing.codeParseFull}
\difficulty{3}
For every bit string $x$, the result of the machine parser on $x$ equals the
result of the suffix-retaining machine parser on $x$ with the suffix component
discarded: both fail together, and on success they return the same machine.
\end{lemma}

\begin{lemma}[Scanner as projection of the suffix-retaining parser]
\lean{Turing.codeScan_full}
\label{Turing.codeScan_full}
\leanok
\uses{Turing.Action, Turing.CodeTM, Turing.codeBitsNat, Turing.codeEraseAction, Turing.codeEraseFin, Turing.codeEraseSymbols, Turing.codeEraseVec, Turing.codeErase_bind, Turing.codeParseFull, Turing.codeReadAction, Turing.codeReadSymbols, Turing.codeScan}
\difficulty{5}
For every bit string $x$, the result of the erased suffix scanner on $x$ equals
the result of the suffix-retaining machine parser on $x$ with the machine
component discarded: both fail together, and on success the scanner returns
exactly the unconsumed suffix that the parser would return.
\end{lemma}

\begin{lemma}[Soundness of the suffix-retaining parser]
\lean{Turing.codeParseFull_sound}
\label{Turing.codeParseFull_sound}
\leanok
\uses{Turing.CodeTM, Turing.CodeTM.serialize, Turing.codeBitsNat, Turing.codePairDecode_sound, Turing.codeParseFull, Turing.codeReadAction_sound, Turing.codeReadFin_sound, Turing.codeReadSymbols_sound, Turing.codeReadVec_sound, Turing.pairEncode}
\difficulty{5}
If the suffix-retaining machine parser succeeds on a bit string $x$, returning
a coded machine $M$ and a suffix $t$, then $x$ is exactly the canonical
serialization of $M$ followed by $t$.
\end{lemma}

\begin{definition}[Canonical serialization of a string]
\lean{Turing.codeCanonical}
\label{Turing.codeCanonical}
\leanok
\uses{Turing.CodeTM.serialize, Turing.codeFallback, Turing.codeScan}
The canonical serialization denoted by a bit string $x$: if the erased suffix
scanner accepts $x$ with unconsumed suffix $t$, the result is the prefix of $x$
of length $|x| - |t|$ (the consumed part); if the scanner rejects $x$, the
result is the serialization of the fixed fallback machine.
\end{definition}

\begin{lemma}[Canonizer computes the decoded serialization]
\lean{Turing.codeCanonical_eq}
\label{Turing.codeCanonical_eq}
\leanok
\uses{Turing.CodeTM.serialize, Turing.codeCanonical, Turing.codeDecode, Turing.codeParseFull, Turing.codeParseFull_sound, Turing.codeParse_full, Turing.codeScan_full}
\difficulty{3}
For every bit string $x$, the canonical serialization denoted by $x$ equals the
serialization of the coded machine that $x$ decodes to under the total decoder.
\end{lemma}
